\documentclass[11pt]{report}
\usepackage[margin=1in]{geometry}
\usepackage{amsmath,amssymb,amsthm}
\usepackage{graphicx}
\usepackage{booktabs}
\usepackage{xcolor}
\usepackage{hyperref}
\usepackage{fancyhdr}
\usepackage{listings}
\usepackage{titlesec}
\usepackage{enumitem}
\usepackage{array}

\definecolor{navy}{RGB}{0,33,71}
\definecolor{steel}{RGB}{70,90,110}
\definecolor{codebg}{RGB}{245,245,248}

\titleformat{\chapter}[display]{\normalfont\bfseries\color{navy}}{\chaptertitlename\ \thechapter}{10pt}{\Large}
\titleformat{\section}{\normalfont\Large\bfseries\color{navy}}{\thesection}{1em}{}
\titleformat{\subsection}{\normalfont\large\bfseries\color{steel}}{\thesubsection}{1em}{}

\lstset{
  backgroundcolor=\color{codebg},
  basicstyle=\ttfamily\small,
  frame=single,
  breaklines=true,
  columns=fullflexible,
  language=Python
}

\pagestyle{fancy}
\fancyhf{}
\rhead{\small\color{steel}AI/ML Modeler — Liquid Financing}
\lhead{\small\color{steel}Barclays Take-Home Research Portfolio}
\rfoot{\thepage}

\title{
  \vspace{-1.5cm}
  {\Large\color{steel}A Dissertation-Style Research Presentation}\\[0.4cm]
  {\huge\bfseries\color{navy} Machine Learning and Generative AI Systems\\ for Cross-Asset Liquid Financing}\\[0.6cm]
  {\large\color{steel}Five Production-Grade Case Studies: Securities-Lending Fee Forecasting,\\
  Client Margin Optimization, Funding-Spread Anomaly Detection,\\
  Deep-Learning Balance Forecasting, and a Retrieval-Augmented Financing Copilot}
}
\author{Candidate Presentation \\ \small Prepared for: Barclays Quantitative Research \& Trading Panel (5 Reviewers)}
\date{Take-Home Discussion \\ Wednesday, July 8, 2026}

\begin{document}
\maketitle
\thispagestyle{empty}

\begin{abstract}
\noindent This work presents five production-oriented machine learning and generative AI systems
designed for a cross-asset Liquid Financing platform spanning Equities \& Delta One, Rate and Credit
Financing, FX, Risk, and Futures \& Prime Derivatives Solutions. Each system is scoped directly against
a stated business roadmap of 50+ identified problem statements and mirrors the explicit technical
requirements of the target role: regularized regression (OLS, LASSO, Elastic Net), tree-ensemble
methods, time-series forecasting, deep sequence models (RNN/LSTM/GRU), and generative AI
(fine-tuning, prompt engineering, retrieval-augmented generation, and evaluation). We present, for
each system: (i) the business problem and P\&L or risk lever it addresses; (ii) the mathematical
formulation and model architecture; (iii) a reference implementation; and (iv) a validation and
production-monitoring protocol consistent with institutional model-risk practice. We conclude with a
cross-cutting discussion of deployment standards, model governance, and the trade-offs between
prompt-engineered and fine-tuned generative systems in a regulated financial environment.
\end{abstract}

\tableofcontents

\chapter{Introduction and Problem Framing}

\section{Business Context}
Liquid Financing intermediates the price of \emph{borrowing} — securities, cash, and balance sheet —
across a cross-asset platform. Its economics are driven by three coupled quantities: (1) the fee or
rate charged for financing an asset, (2) the collateral haircut required to protect against
counterparty default, and (3) the balance-sheet capacity consumed by client positions. Each of the
five case studies in this dissertation targets exactly one of these economic levers, chosen so that,
collectively, they cover every technical bullet point in the target job description while remaining
traceable to a concrete desk P\&L or risk outcome.

\section{Mapping to the Stated Technical Requirements}
\begin{center}
\begin{tabular}{p{2.7cm}p{6.3cm}p{5.5cm}}
\toprule
\textbf{Case Study} & \textbf{Core Method} & \textbf{JD Requirement Addressed} \\
\midrule
P1 & Elastic Net + Gradient-Boosted Trees + AR/event time-series blend & Regression (OLS/LASSO/Elastic Net), Tree models, Time series \\
P2 & LASSO haircut model + GBM correlation-breakdown add-on & Regression, Tree models, model-risk governance \\
P3 & Robust Mahalanobis + Isolation Forest + LLM text fusion & Time series, Gen AI \\
P4 & Sequence-to-sequence GRU/LSTM with quantile heads & Deep Learning (Neural Networks, RNN, LSTM, GRU) \\
P5 & Hybrid retrieval-augmented LLM copilot & Gen AI: fine-tuning, prompting, RAG, evaluation, inference infrastructure \\
\bottomrule
\end{tabular}
\end{center}

\chapter{P1 --- Securities-Lending Fee and Rebate-Rate Forecasting}

\section{Problem Statement}
We forecast the daily specialness fee $f_{i,t}$, expressed in annualized basis points over the general
collateral (GC) rate, for a universe of hard-to-borrow equities, at horizons $h \in \{1,\dots,5\}$
business days, to feed a pre-trade financing pricing engine.

\section{Model Formulation}
A three-tier ensemble is used. The first tier is an Elastic Net regression:
\begin{equation}
\hat\beta = \arg\min_{\beta}\; \lVert y - X\beta \rVert_2^2 + \lambda_1 \lVert \beta \rVert_1 + \lambda_2 \lVert \beta \rVert_2^2
\end{equation}
chosen for auditability: every surviving coefficient can be defended in a model-risk review.
The second tier is a gradient-boosted tree ensemble,
\begin{equation}
F_M(x) = \sum_{m=1}^{M} \gamma_m h_m(x), \qquad h_m = \arg\min_h \sum_{i} L\big(y_i,\, F_{m-1}(x_i) + h(x_i)\big),
\end{equation}
which captures the empirically convex relationship between utilization and fee once utilization
exceeds approximately 90\%, a threshold effect that a purely linear model cannot represent. The third
tier is a mean-reverting, event-augmented time-series component,
\begin{equation}
f_{i,t} = \mu_i + \phi\,(f_{i,t-1} - \mu_i) + \sum_j \kappa_j \, \mathbb{1}[\text{event}_{j,t}]\, e^{-\eta_j (t - t_j)} + \varepsilon_t,
\end{equation}
which explicitly models the predictable fee spikes around dividend record dates and index
rebalancing events. The final forecast is a constrained convex combination of the three tiers, with
weights fit by non-negative least squares on a walk-forward validation window.

\section{Reference Implementation (excerpt)}
\begin{lstlisting}
en  = ElasticNetCV(l1_ratio=[.1,.5,.7,.9,.95,1], cv=5).fit(X_train, y_train)
gbm = LGBMRegressor(n_estimators=600, num_leaves=31,
                     learning_rate=0.03).fit(X_train, y_train)
w, _ = nnls(np.column_stack([en.predict(X_val), gbm.predict(X_val)]), y_val)
forecast = np.column_stack([en.predict(X_new), gbm.predict(X_new)]) @ (w / w.sum())
\end{lstlisting}

\section{Validation Protocol}
Weighted MAPE (notional-weighted), directional hit-rate on fee changes exceeding 25 bps, and pinball
loss at the 10th/90th percentiles are tracked on an expanding walk-forward window. Feature population
stability index (PSI) is monitored monthly, with automatic fallback to the Elastic Net tier if the
gradient-boosted tree's input PSI exceeds 0.25.

\chapter{P2 --- Client Margin and Haircut Optimization}

\section{Problem Statement}
Given a client's cross-asset collateral basket, we predict the required per-asset haircut and
portfolio-level initial margin such that the modeled 99\% five-day liquidation shortfall is covered,
while minimizing over-collateralization relative to the existing rules-based haircut grid.

\section{Model Formulation}
The per-asset haircut is modeled log-linearly with LASSO-selected covariates,
\begin{equation}
\log h_i = \beta_0 + \beta_1 \log \sigma_i + \beta_2 \log(\text{ADV}_i) + \beta_3\, \text{Corr}_{i,\text{mkt}} + \beta_4\, \text{Rating}_i + \varepsilon_i,
\end{equation}
and the portfolio-level required initial margin is decomposed into a linear per-asset component plus
a tree-learned correlation-breakdown add-on,
\begin{equation}
\text{IM}_{\text{portfolio}} = \sum_i h_i\,|q_i|\,P_i \;+\; g_\theta(\mathbf{q}, \boldsymbol\Sigma, \text{concentration}),
\end{equation}
where $g_\theta$ is a gradient-boosted regressor trained on the historical-simulation shortfall
residual unexplained by the linear component. This add-on is precisely the mechanism that captures
correlations moving toward one under stress, historically the primary failure mode of purely
rules-based haircut grids. The final output is clipped to $[\text{rules-based floor},\, 1.2\times
\text{VaR}_{99\%}]$ and trained against the chance constraint
\begin{equation}
\Pr\big(\text{Shortfall}_{5d} > h \cdot V\big) \le 1\%.
\end{equation}

\section{Backtesting Discipline}
Kupiec proportion-of-failures and Christoffersen independence tests are applied to daily VaR breach
counts. Champion/challenger comparison runs for a minimum of two full quarters in shadow mode before
promotion, with stress overlays replaying 2008, March 2020, and 2023 regional-banking scenarios
against the current book.

\chapter{P3 --- Cross-Asset Funding-Spread Anomaly Detection}

\section{Problem Statement}
We detect anomalous divergence across the repo, securities-lending, cross-currency-basis, and
credit-basis complex in real time, fused with a natural-language classifier over overnight desk
commentary, and rank alerts by likely P\&L impact.

\section{Model Formulation}
A robust multivariate anomaly score is computed via Mahalanobis distance under an EWMA-updated,
Minimum-Covariance-Determinant-robustified covariance estimate,
\begin{equation}
z_t = (\mathbf{x}_t - \boldsymbol\mu_t)^\top \boldsymbol\Sigma_t^{-1} (\mathbf{x}_t - \boldsymbol\mu_t),
\end{equation}
cross-checked by an Isolation Forest score $s(x) = 2^{-E[h(x)]/c(n)}$ that is more robust to
regime-dependent non-Gaussianity. A fine-tuned or prompt-engineered language model classifies
overnight commentary into stress categories, and the statistical and textual signals are fused as
\begin{equation}
\text{AlertScore}_t = \alpha\cdot \text{sigmoid}(z_t - z^*) + (1-\alpha)\cdot p(\text{stress} \mid \text{text}_t).
\end{equation}

\section{Human-in-the-Loop Design}
Every alert above threshold routes to a trader queue accompanied by the driving spread chart, the
top-attributed statistical features, and a cited natural-language rationale; no action is executed
automatically. Precision and recall are tracked weekly against trader-labeled outcomes, with the
alert threshold retuned via a Neyman--Pearson-style asymmetric cost ratio reflecting that missed
funding-stress events are materially more costly than false positives on this desk.

\chapter{P4 --- Prime Balance and Utilization Forecasting via Deep Sequence Models}

\section{Problem Statement}
We forecast daily aggregate client financing balances and utilization ratios one to twenty business
days ahead, with explicit treatment of month-end and quarter-end window-dressing effects, for
Treasury capacity and RWA/LCR planning.

\section{Model Formulation}
A sequence-to-sequence Gated Recurrent Unit (GRU) encoder-decoder is used,
\begin{equation}
h_t = \text{GRU}(x_t, h_{t-1}), \qquad \hat y_{t+1:t+H} = \text{Decoder}(h_T,\; \text{calendar embeddings}),
\end{equation}
with an explicit, learned calendar embedding concatenated at every decoder step so that the
known, calendar-driven quarter-end balance-reduction pattern is treated as an inspectable input rather
than left for the network to rediscover implicitly. Three quantile heads are trained jointly via the
pinball loss,
\begin{equation}
\mathcal{L} = \sum_{q\in\{0.1,0.5,0.9\}} \sum_{h=1}^{H} \rho_q\big(y_{t+h} - \hat y_{t+h}^{(q)}\big), \qquad \rho_q(u) = u\,(q - \mathbb{1}[u<0]),
\end{equation}
producing a P10/P50/P90 balance forecast band rather than a false-precision point estimate. An LSTM
variant is retained as a challenger given its stronger long-memory gating.

\section{Baseline Discipline}
The GRU architecture is required to outperform a seasonal-naive baseline, a SARIMAX model with an
explicit quarter-end dummy, and a well-tuned LightGBM model on lagged and calendar features, on
walk-forward mean absolute percentage error and pinball loss, before its added production and serving
complexity is justified to the panel.

\chapter{P5 --- A Retrieval-Augmented Generation Copilot for the Financing Desk}

\section{Problem Statement}
We build a retrieval-augmented generation (RAG) system over internal financing policy documents, desk
commentary, and the structured outputs of the P1--P4 models, that answers trader queries with
citations and explicitly declines to answer when evidence is insufficient.

\section{Retrieval Formulation}
Hybrid sparse-dense retrieval is used,
\begin{equation}
\text{score}(q,d) = \lambda \cdot \text{BM25}(q,d) + (1-\lambda)\cdot \cos\big(E(q), E(d)\big),
\end{equation}
followed by cross-encoder reranking of the top candidates, with $\lambda$ tuned against a
human-labeled query/document relevance set using recall@10 as the primary retrieval metric.

\section{Fine-Tuning versus Prompt Engineering: An Explicit Trade-off}
Per the target role's requirement to \emph{evaluate and articulate trade-offs} between modeling
approaches, we adopt a staged strategy: a prompt-engineered, retrieval-augmented system is deployed
first to establish an evaluation harness and collect production query logs; intent classification and
domain-jargon normalization are subsequently migrated to a small, LoRA-fine-tuned model once a
sufficient labeled query corpus exists, while final-answer synthesis remains on a frontier base model
with retrieval for auditability. This reflects a governance-first calculus: full fine-tuning offers
superior domain fluency and lower marginal cost per query, but at the expense of update latency and
audit-trail complexity relative to a prompt-plus-retrieval design whose full reasoning trace is
visible at inference time.

\section{Evaluation Harness}
\begin{equation}
\text{Faithfulness} = \frac{\#\{\text{claims in the answer supported by retrieved context}\}}{\#\{\text{claims in the answer}\}}.
\end{equation}
Faithfulness, retrieval recall@10, human-graded answer correctness against a quarterly-refreshed
golden set, and refusal precision on out-of-scope queries are tracked continuously, with a p95
latency service-level objective of under four seconds during trading hours.

\section{Inference Infrastructure}
Consistent with the target role's requirement to configure inference infrastructure, non-public,
client-identifying, or desk-proprietary queries are routed to an on-premises, BMC-hosted open-weight
model behind a request-routing layer that enforces entitlement and PII filtering prior to prompt
assembly; only public-market-context queries are eligible for routing to an external frontier model
endpoint.

\chapter{Cross-Cutting Production and Governance Standards}

\section{Validation}
All five systems are validated exclusively on walk-forward, expanding-window splits; random
cross-validation is explicitly avoided given the autocorrelated, event-clustered structure of
financing time series, which would otherwise leak information across the train/validation boundary.

\section{Model Risk Documentation}
Each system carries a standard model-risk package: statement of purpose, conceptual soundness review,
developmental evidence (in-sample and out-of-sample performance), independent outcomes analysis
(backtesting), and an ongoing performance-monitoring plan, with a minimum two-quarter shadow period
prior to any production promotion.

\section{Monitoring and Fallback}
Feature-level population stability index, prediction drift, and backtest breach counts are monitored
continuously, with automatic fallback to the prior champion model on threshold breach, and canary
release to a 5\% traffic slice prior to full production rollout.

\chapter{Conclusion}
The five systems presented — fee forecasting, margin optimization, anomaly detection, deep
sequence-model balance forecasting, and a retrieval-augmented financing copilot — jointly span every
modeling technique enumerated in the target job description, are each traceable to a specific desk
P\&L or risk lever, and are each specified with a production validation and governance protocol
suitable for a regulated cross-asset financing business. Collectively they represent a candidate
execution plan against the stated backlog of fifty-plus prioritized problem statements, beginning with
the ten already flagged as active priorities.

\end{document}
